% Part: computability
% Chapter: computability-theory
% Section: russells-paradox

\documentclass[../../../include/open-logic-section]{subfiles}

\begin{document}

\olfileid{cmp}{thy}{rus}
\olsection{مقارنة بمفارقة راسل}

ويتضح شأن هذه الحجج بموازنتها بمفارقة راسل، لنرى ما يأتلف منها وما
يفترق:
\begin{enumerate}
\item مفارقة راسل: لتكن $\OLId{latin-looped}{S}=\Setabs{\OLId{latin-ordinary}{x}}{\OLId{latin-ordinary}{x}\notin \OLId{latin-ordinary}{x}}$. عندئذ
$\OLId{latin-looped}{S}\in \OLId{latin-looped}{S}$ إذا وفقط إذا كان $\OLId{latin-looped}{S}\notin \OLId{latin-looped}{S}$، وهذا تناقض.

  \emph{النتيجة:} يمتنع وجود المجموعة~$\OLId{latin-looped}{S}$ على هذه الصفة؛ فإثبات «مجموعة جميع
  المجموعات» يناقض سائر بديهيات نظرية المجموعات.

\item تعديل لمفارقة راسل: لتكن $\OLId{latin-looped}{F}$ «الدالة» من مجموعة جميع الدوال إلى
$\{\OLMathNumeral{0},\OLMathNumeral{1}\}$، المعرفة بـ
  \[
  \OLId{latin-looped}{F}(\OLId{latin-ordinary}{f}) =
  \begin{cases}
    \OLMathNumeral{1} & \text{إذا كانت $f$ في مجال $f$ وكان $f(f) = 0$} \\
    \OLMathNumeral{0} & \text{في غير ذلك}
  \end{cases}
  \]
  وبمثل الحجة الأولى يلزم أن $\OLId{latin-looped}{F}(\OLId{latin-looped}{F})=\OLMathNumeral{0}$ إذا وفقط إذا كان $\OLId{latin-looped}{F}(\OLId{latin-looped}{F})=\OLMathNumeral{1}$، وهذا تناقض.

  \emph{النتيجة:} لا تكون $\OLId{latin-looped}{F}$ دالة؛ إذ «مجموعة جميع الدوال» أعظم من أن
  يصح اتخاذها مجالًا لدالة.

\item الحجة القطرية: لتكن $\OLId{latin-ordinary}{f}_\OLMathNumeral{0}$ و$\OLId{latin-ordinary}{f}_\OLMathNumeral{1}$ و… تعداد الدوال الجزئية
  القابلة للحساب، ولتكن $\OLId{latin-looped}{G}\colon\Nat\to\{\OLMathNumeral{0},\OLMathNumeral{1}\}$ معرفة بـ
  \[
  \OLId{latin-looped}{G}(\OLId{latin-ordinary}{x}) =
  \begin{cases}
    \OLMathNumeral{1} & \text{إذا كان $f_x(x)\downarrow = 0$} \\
    \OLMathNumeral{0} & \text{في غير ذلك}
  \end{cases}
  \]
  فلو كانت $\OLId{latin-looped}{G}$ قابلة للحساب، لكانت $\OLId{latin-ordinary}{f}_\OLId{latin-ordinary}{k}$ لعدد ما $\OLId{latin-ordinary}{k}$، وحينئذ
  يكون $\OLId{latin-looped}{G}(\OLId{latin-ordinary}{k})=\OLMathNumeral{1}$ إذا وفقط إذا كان $\OLId{latin-looped}{G}(\OLId{latin-ordinary}{k})=\OLMathNumeral{0}$، وهذا تناقض.

  \emph{النتيجة:} تمتنع قابلية $\OLId{latin-looped}{G}$ للحساب، ولا يمتنع أن تكون $\OLId{latin-looped}{G}$ دالة في
  نظرية المجموعات. فليس هذا تناقضًا في البديهيات، وإنما هو تمييز لما يقبل الحساب مما لا يقبله.
\end{enumerate}

وقد يلتبس تداخل هذه الأسماء: الدوال الجزئية، والدوال القابلة للحساب،
والدوال الجزئية القابلة للحساب. فلنميزها. الدوال
الجزئية من $\Nat$ إلى $\Nat$ تؤلف مجموعة واسعة؛ منها الكلي،
ومنها القابل للحساب، ومنها ما يجمع الأمرين، ومنها ما يفقدهما
معًا. واصطلاحنا، حين نطلق «دالة» بلا قيد، أن نريد الكلية.
فتكون الأصناف كما يأتي:
\begin{enumerate}
\item دوال قابلة للحساب؛
\item دوال جزئية قابلة للحساب وليست كلية؛
\item دوال غير قابلة للحساب؛
\item دوال جزئية ليست كلية ولا قابلة للحساب.
\end{enumerate}
ولتستحضر هذه الأصناف، ارسم مربعًا يمثل جميع الدوال الجزئية من
$\Nat$ إلى $\Nat$، واجعل فيه منطقتين متداخلتين: إحداهما للدوال الكلية،
والأخرى للجزئية القابلة للحساب. ثم تمرّن بأن
تصف دالة في كل منطقة ينقسم إليها الرسم.

\end{document}
